% !TEX program = lualatex
\documentclass[a4paper,11pt]{ltjsarticle}

\usepackage{amsmath,amssymb}
\usepackage{booktabs}
\usepackage{geometry}
\usepackage{graphicx}
\usepackage{hyperref}
\geometry{margin=24mm}

\title{Seminar 用本実験結果メモ\\
\large Friedman-I における \texttt{bo\_current} の評価}
\author{}
\date{2026年6月2日}

\newcommand{\HV}{\mathrm{HV}}

\begin{document}
\maketitle

\section{目的}

本メモは，Seminar 発表に向けた本実験 for Seminar として，
現時点で採用する方針である \texttt{bo\_current} を
Friedman-I シンボリック回帰問題に適用した結果をまとめるものである．

本研究の提案法は，多目的 GP の世代状態を観測し，
文脈付き BO により交叉率 \(p_c\)，突然変異率 \(p_m\)，
更新周期 \(k\) を逐次決定する閉ループ制御手法である．
今回の Seminar 用実験では，後続検討候補である
\(\beta_{\mathrm{EI}}\) ではなく，
既定の \texttt{bo\_current} 方針を本線として用いる．

\section{対象問題}

対象問題は Friedman-I 型のシンボリック回帰である．
入力変数を \(x_1,\ldots,x_5\) とすると，目標関数は
\[
  y =
  10\sin(\pi x_1 x_2)
  +20(x_3-0.5)^2
  +10x_4
  +5x_5
\]
で与えられる．
各入力は一様分布 \(x_i\sim U(0,1)\) から生成し，
訓練点数は 200，テスト点数は 200 とした．

GP の目的関数は 2 目的である．
第1目的は訓練データに対する正規化 RMSE
（training NRMSE）の最小化であり，
第2目的は式木サイズの最小化である．
したがって本実験では，精度の良さと式の簡潔さのトレードオフを
Pareto front として評価する．

\section{比較条件}

比較対象は，固定率 GP 3 条件と，提案法 \texttt{bo\_current} である．
全手法で同じ多目的 GP エンジンを用い，違いは
交叉率・突然変異率・更新周期の決定方法のみに限定した．

\begin{table}[htbp]
\centering
\caption{比較した制御条件}
\begin{tabular}{lll}
\toprule
手法 & 設定 & 位置づけ \\
\midrule
\texttt{plain\_fixed\_standard}
& \(p_c=0.80,\ p_m=0.05\)
& 標準的な固定率 GP \\
\texttt{plain\_fixed\_high\_mutation}
& \(p_c=0.70,\ p_m=0.20\)
& 突然変異を強めた固定率 GP \\
\texttt{plain\_fixed\_high\_crossover}
& \(p_c=0.90,\ p_m=0.05\)
& 交叉を強めた固定率 GP \\
\texttt{bogp\_current}
& \(p_c,p_m,k\) を BO で逐次決定
& 提案法の現行版 \\
\bottomrule
\end{tabular}
\end{table}

\texttt{bo\_current} では，
\[
  k\in\{1,3,5\}
\]
を候補とし，\(k\) ごとに GP surrogate を持つ．
EI の改善基準値は各 \(k\) 内の過去最高報酬であり，
これは \texttt{per-k} EI と解釈できる．
すなわち，今回の \texttt{bo\_current} は
\[
  r_{\mathrm{best}}^{(k)}
  =
  \max\{r_i\mid k_i=k\}
\]
を用いて EI を計算する．

\section{実験条件}

実験条件を表\ref{tab:experiment-setting}に示す．
各手法は seed 0 から 99 までの 100 試行で評価した．

\begin{table}[htbp]
\centering
\caption{Seminar 用本実験の条件}
\label{tab:experiment-setting}
\begin{tabular}{ll}
\toprule
項目 & 設定 \\
\midrule
対象問題 & Friedman-I symbolic regression \\
seed & \(0,\ldots,99\) \\
集団サイズ & 24 \\
世代数 & 30 \\
評価回数 & 720 / run \\
目的数 & 2 \\
目的関数 & training NRMSE，tree size \\
archive key & topology + node value \\
HV 参照点 & \((1.0,1.0)\) \\
\bottomrule
\end{tabular}
\end{table}

\section{結果}

\subsection{最終 HV と多様性}

表\ref{tab:final-results}に，100 seed における final archive HV，
final diversity，archive size の平均と標準偏差を示す．

\begin{table}[htbp]
\centering
\caption{Friedman-I における最終指標}
\label{tab:final-results}
\resizebox{\linewidth}{!}{%
\begin{tabular}{lrrrrr}
\toprule
手法 & HV 平均 & HV 標準偏差 & 多様性平均 & 多様性標準偏差 & archive size 平均 \\
\midrule
\texttt{plain\_fixed\_standard}
& 0.757001 & 0.076593 & 0.601648 & 0.126511 & 19.15 \\
\texttt{plain\_fixed\_high\_mutation}
& 0.776600 & 0.051848 & 0.609492 & 0.113961 & 14.00 \\
\texttt{plain\_fixed\_high\_crossover}
& 0.764665 & 0.072106 & 0.615219 & 0.102404 & 19.30 \\
\texttt{bogp\_current}
& 0.773081 & 0.055720 & 0.603050 & 0.123937 & 15.72 \\
\bottomrule
\end{tabular}
}
\end{table}

図\ref{fig:final-hv-diversity}に，final archive HV と final diversity を
平均 \(\pm\) 標準偏差で示す．
\texttt{bogp\_current} は標準固定率 GP および高交叉固定率 GP よりも
HV 平均が高い．
一方で，高突然変異固定率 GP が HV 平均で最も高く，
\texttt{bogp\_current} はそれにわずかに届いていない．

\begin{figure}[htbp]
  \centering
  \includegraphics[width=\linewidth]{outputs/seminar_bo_current/sr_alpha_friedman/seminar_bo_current_seed100_20260602/seminar_final_hv_diversity.png}
  \caption{final archive HV と final diversity の比較}
  \label{fig:final-hv-diversity}
\end{figure}

\subsection{seed 内差分}

表\ref{tab:paired-diff}に，同じ seed 内で
\texttt{bogp\_current} と各固定率 GP の final HV 差分を取った結果を示す．
差分は
\[
  \Delta \HV =
  \HV_{\texttt{bogp\_current}}
  -
  \HV_{\texttt{baseline}}
\]
として定義した．

\begin{table}[htbp]
\centering
\caption{\texttt{bogp\_current} と固定率 GP の seed 内 final HV 差分}
\label{tab:paired-diff}
\begin{tabular}{lrrrrr}
\toprule
比較 & 差分平均 & 差分標準偏差 & win & tie & loss \\
\midrule
\texttt{bogp\_current} - standard
& 0.016080 & 0.089679 & 64 & 4 & 32 \\
\texttt{bogp\_current} - high mutation
& -0.003519 & 0.068742 & 40 & 4 & 56 \\
\texttt{bogp\_current} - high crossover
& 0.008416 & 0.088748 & 54 & 2 & 44 \\
\texttt{bogp\_current} - best fixed per seed
& -0.020092 & 0.053069 & 21 & 1 & 78 \\
\bottomrule
\end{tabular}
\end{table}

図\ref{fig:paired-diff}は，表\ref{tab:paired-diff}の差分を図示したものである．
\texttt{bogp\_current} は標準固定率 GP に対して平均差が正であり，
100 seed 中 64 seed で勝っている．
しかし，高突然変異固定率 GP に対しては平均差が負であり，
best fixed per seed と比較すると 21 seed でのみ勝っている．

\begin{figure}[htbp]
  \centering
  \includegraphics[width=0.92\linewidth]{outputs/seminar_bo_current/sr_alpha_friedman/seminar_bo_current_seed100_20260602/seminar_paired_hv_difference.png}
  \caption{\texttt{bogp\_current} と固定率 GP の seed 内 final HV 差分}
  \label{fig:paired-diff}
\end{figure}

\subsection{HV 分布}

図\ref{fig:hv-boxplot}に，100 seed における final archive HV の分布を示す．
中央値を見ると，各手法の差は大きくない．
一方で，標準固定率 GP と高交叉固定率 GP では低 HV 側の外れ値が見られる．
\texttt{bogp\_current} も外れ値を完全には避けられていないが，
分布の中心は標準固定率 GP より高い位置にある．

\begin{figure}[htbp]
  \centering
  \includegraphics[width=0.92\linewidth]{outputs/seminar_bo_current/sr_alpha_friedman/seminar_bo_current_seed100_20260602/seminar_final_hv_boxplot.png}
  \caption{100 seed における final archive HV の分布}
  \label{fig:hv-boxplot}
\end{figure}

\subsection{\texorpdfstring{\(k\)}{k} の選択傾向}

\texttt{bogp\_current} における更新周期 \(k\) の選択回数を図\ref{fig:k-counts}に示す．
100 seed 全体での選択回数は
\[
  k=1:491,\quad
  k=3:333,\quad
  k=5:302
\]
であった．
平均制御更新回数は 11.26 回であり，毎世代更新の 30 回より少ない．
これは，\texttt{bo\_current} が毎世代細かく制御するのではなく，
必要に応じて数世代単位で制御入力を保持する挙動を示している．

\begin{figure}[htbp]
  \centering
  \includegraphics[width=0.62\linewidth]{outputs/seminar_bo_current/sr_alpha_friedman/seminar_bo_current_seed100_20260602/seminar_k_selection_counts.png}
  \caption{\texttt{bogp\_current} における更新周期 \(k\) の選択回数}
  \label{fig:k-counts}
\end{figure}

\section{考察}

今回の Friedman-I 実験から，\texttt{bo\_current} は
標準固定率 GP に対して final archive HV を改善する傾向を示した．
これは，GP の世代状態を観測して \(p_c,p_m,k\) を切り替える
閉ループ制御の枠組みが，少なくとも標準的な固定率設定の単純な置き換えとして
動作していることを示す結果である．

一方で，高突然変異固定率 GP が HV 平均で最良であったことは重要である．
この結果は，現行の \texttt{bo\_current} が
「十分に調整された固定率 GP」を安定して上回る段階にはまだ達していないことを示している．
したがって，Seminar では
「提案法が全 baseline に勝った」と主張するのではなく，
次のように位置づけるのが妥当である．

\begin{itemize}
  \item 提案法の閉ループ制御は実装上動作しており，標準固定率 GP より良い傾向を示した．
  \item しかし，高突然変異固定率 GP にはわずかに届かず，報酬設計や BO 制御器の改善余地が残る．
  \item \(k\) を制御対象に含めることで，毎世代更新ではなく平均 11.26 回の制御更新で進化を進められた．
  \item 今後は，固定率のグリッド比較，報酬内訳分析，\(k\) ごとの報酬分布分析を行い，制御方策を改良する必要がある．
\end{itemize}

\section{Seminar 発表での説明方針}

Seminar では，結果を過度に強く主張するよりも，
「提案法の枠組みが動いたこと」と
「現行版の課題が明確になったこと」を主な成果として説明する．
特に，今回の結果は次のストーリーで説明しやすい．

\begin{enumerate}
  \item 固定率 GP では，どの \(p_c,p_m\) が良いかは問題依存である．
  \item Friedman-I では，高突然変異固定率が強い baseline になった．
  \item \texttt{bo\_current} は標準固定率より HV を改善したが，最良固定率には届かなかった．
  \item これは，閉ループ制御の方向性は妥当だが，報酬設計や BO の比較基準をさらに詰める必要があることを示す．
  \item Seminar 後には，\(\beta_{\mathrm{EI}}\) などの EI 改良案，および \(k\) の扱いを再検討する．
\end{enumerate}

\end{document}
